% !Mode:: "TeX:UTF-8"
% !TEX root = ../../main.tex
% ============================================================
% TOPIC 20
% ============================================================
\mashgulot{20}{Linear operators with a basis of eigenvectors}

\begin{teorema}
The matrix of an operator $\mathcal A$ on an $n$-dimensional space is
diagonal in some basis if and only if $\mathcal A$ has $n$ linearly
independent eigenvectors, that is, if and only if its eigenvectors form
a basis. In that case
\[\kf{T^{-1}AT=D=\begin{pmatrix}\lambda_1&&0\\&\ddots&\\0&&\lambda_n\end{pmatrix},}\]
where the columns of $T$ are the eigenvectors and the diagonal of $D$
carries the corresponding eigenvalues.
\end{teorema}

\textbf{Sufficient conditions.} If the $n$ eigenvalues of $A$ are
pairwise distinct, then $A$ is diagonalizable. Every real symmetric
matrix is diagonalizable, and even by an orthogonal matrix $Q$:
$Q^TAQ=D$. When $A$ has a multiple eigenvalue $\lambda$, the matrix is
diagonalizable if and only if, for every such $\lambda$, the
\emph{geometric} multiplicity
$\dim\ker(A-\lambda E)=n-\operatorname{rank}(A-\lambda E)$ equals the
\emph{algebraic} multiplicity of $\lambda$ as a root of the
characteristic polynomial.

\begin{misol}
Diagonalize $A=\begin{pmatrix}4&1\\2&3\end{pmatrix}$ and find $A^n$.

\textbf{Solution.} In the previous section we found $\lambda_1=2$ with
eigenvector $(1;-2)$ and $\lambda_2=5$ with eigenvector $(1;1)$.
Therefore
\[\kf{T=\begin{pmatrix}1&1\\-2&1\end{pmatrix},\quad
T^{-1}=\frac13\begin{pmatrix}1&-1\\2&1\end{pmatrix},\quad
T^{-1}AT=\begin{pmatrix}2&0\\0&5\end{pmatrix}.}\]
Since $A=TDT^{-1}$, we have $A^n=TD^nT^{-1}$, and multiplying out gives
\[\kf{A^n=\frac13\begin{pmatrix}2^n+2\cdot5^n & 5^n-2^n\\ 2\cdot5^n-2^{n+1} & 2^{n+1}+5^n\end{pmatrix}.}\]
For instance, for $n=2$ we obtain
$A^2=\frac13\begin{pmatrix}54&21\\42&33\end{pmatrix}=\begin{pmatrix}18&7\\14&11\end{pmatrix}$,
which agrees with direct multiplication.
\end{misol}

\begin{misol}
Is $A=\begin{pmatrix}2&1\\0&2\end{pmatrix}$ diagonalizable?

\textbf{Solution.} The only eigenvalue is $\lambda=2$, of algebraic
multiplicity two. The matrix $A-2E=\begin{pmatrix}0&1\\0&0\end{pmatrix}$
has rank 1, so the eigenspace is one-dimensional and is spanned by
$(1;0)$. There is no pair of linearly independent eigenvectors, and
$A$ is \textbf{not diagonalizable}; it is a Jordan block.
\end{misol}

\begin{misol}
Diagonalize $A=\begin{pmatrix}3&1\\1&3\end{pmatrix}$ by an orthogonal
matrix.

\textbf{Solution.} The characteristic equation $\lambda^2-6\lambda+8=0$
gives $\lambda_1=2$ with eigenvector $(1;-1)$ and $\lambda_2=4$ with
eigenvector $(1;1)$. These vectors are orthogonal; normalizing them, we
obtain
\[\kf{Q=\dfrac{1}{\sqrt2}\begin{pmatrix}1&1\\-1&1\end{pmatrix},\qquad
Q^TAQ=\begin{pmatrix}2&0\\0&4\end{pmatrix}.}\]
\end{misol}

\subsection{Applications of diagonalization in information security}

\textbf{Fast computation of matrix powers.} Many problems in cryptography
and modelling require $A^k$ for large $k$. For example, the state of a
linear recurrent generator (such as an LFSR) is $s_{k}=A^ks_0$. If $A$ is
diagonalizable, then $A^k=TD^kT^{-1}$, and the work reduces to raising
$n$ numbers to the $k$th power.

\begin{misol}
A Fibonacci-type generator is described by
$\binom{F_{k+1}}{F_k}=\begin{pmatrix}1&1\\1&0\end{pmatrix}^k\binom{1}{0}$.
The eigenvalues of this matrix are $\lambda_{1,2}=\frac{1\pm\sqrt5}{2}$,
and diagonalization yields Binet’s formula
$F_k=\dfrac{\lambda_1^k-\lambda_2^k}{\sqrt5}$, which gives the state of
the generator at any step without computing the previous ones.
\end{misol}

\textbf{Markov models.} In a Markov model of user behaviour the state
distribution after $k$ steps is $x_{k}=P^kx_0$. Diagonalization gives the
limiting distribution as $k\to\infty$, and a substantial deviation from
it is treated as an anomaly.

\subsection{Computing with Python, Maple and Mathcad}

We diagonalize the matrix $A=\begin{pmatrix}4&1\\2&3\end{pmatrix}$ of the
first example and use the result to compute $A^{10}$. SymPy’s
\texttt{diagonalize} returns the pair $(T,D)$ directly.
\begin{Verbatim}[fontsize=\footnotesize,frame=single,label={\small Python}]
import sympy as sp
A = sp.Matrix([[4, 1], [2, 3]])
T, D = A.diagonalize()
print(T, D)
print(T.inv() * A * T == D)
print(T * D**10 * T.inv())
\end{Verbatim}
\begin{Verbatim}[fontsize=\footnotesize,frame=single,label={\small Output}]
Matrix([[-1, 1], [2, 1]]) Matrix([[2, 0], [0, 5]])
True
Matrix([[6510758, 3254867], [6509734, 3255891]])
\end{Verbatim}
SymPy chose the eigenvector $(-1;2)$ instead of $(1;-2)$; any nonzero
multiple will do. The entry $6510758=(2^{10}+2\cdot5^{10})/3$ agrees with
the general formula for $A^n$.

In Maple, \texttt{Eigenvectors} returns the eigenvalues together with the
matrix $T$ of eigenvectors, and the zero matrix in the output confirms
that $T^{-1}AT$ is the diagonal matrix built from these eigenvalues.
\begin{Verbatim}[fontsize=\footnotesize,frame=single,label={\small Maple}]
with(LinearAlgebra):
A := Matrix([[4, 1], [2, 3]]):
v, T := Eigenvectors(A):
MatrixInverse(T) . A . T - DiagonalMatrix(v);
                       [0  0]
                       [0  0]
A^10;
               [6510758  3254867]
               [6509734  3255891]
\end{Verbatim}

In Mathcad the matrix of (normalized) eigenvectors is produced by
\texttt{eigenvecs(A)}; the power $A^{10}$ is typed with the usual power
operator. (Mathcad and Maple may list the eigenvalues in a different
order; the columns of $T$ are then permuted accordingly.)
\begin{Verbatim}[fontsize=\footnotesize,frame=single,label={\small Mathcad}]
A := | 4  1 |        V := eigenvecs(A)
     | 2  3 |
V^-1 * A * V = | 2  0 |
               | 0  5 |
A^10 = | 6510758  3254867 |
       | 6509734  3255891 |
\end{Verbatim}

\subsection{Exercises on diagonalization}

\emph{1. Diagonalize the matrix, that is, find $T$ and $D$ such that
$T^{-1}AT=D$.}
\bigskip
\begin{multicols}{2}\raggedright
\begin{enumerate}[label=\arabic*., leftmargin=*]
\item $\begin{pmatrix}-2 & 0\\ 4 & 0\end{pmatrix}$
\item $\begin{pmatrix}-4 & -4\\ 2 & 2\end{pmatrix}$
\item $\begin{pmatrix}7 & -10\\ 5 & -8\end{pmatrix}$
\item $\begin{pmatrix}-2 & 0\\ 1 & -3\end{pmatrix}$
\item $\begin{pmatrix}-6 & -6\\ 3 & 3\end{pmatrix}$
\item $\begin{pmatrix}3 & 0\\ 0 & -3\end{pmatrix}$
\item $\begin{pmatrix}3 & 0\\ 0 & 1\end{pmatrix}$
\item $\begin{pmatrix}-1 & 0\\ -8 & 3\end{pmatrix}$
\item $\begin{pmatrix}4 & 0\\ 0 & 0\end{pmatrix}$
\item $\begin{pmatrix}1 & 0\\ 8 & -3\end{pmatrix}$
\item $\begin{pmatrix}0 & -3\\ 0 & -3\end{pmatrix}$
\item $\begin{pmatrix}0 & 0\\ 0 & 4\end{pmatrix}$
\item $\begin{pmatrix}-1 & 2\\ 0 & -3\end{pmatrix}$
\item $\begin{pmatrix}4 & 12\\ 0 & -2\end{pmatrix}$
\item $\begin{pmatrix}0 & 0\\ -4 & 4\end{pmatrix}$
\item $\begin{pmatrix}-2 & -2\\ 1 & 1\end{pmatrix}$
\item $\begin{pmatrix}4 & 0\\ 0 & -1\end{pmatrix}$
\item $\begin{pmatrix}-2 & 0\\ 0 & -3\end{pmatrix}$
\item $\begin{pmatrix}-5 & 2\\ -1 & -2\end{pmatrix}$
\item $\begin{pmatrix}4 & 0\\ 7 & -3\end{pmatrix}$
\item $\begin{pmatrix}4 & 8\\ -4 & -8\end{pmatrix}$
\item $\begin{pmatrix}0 & 0\\ 8 & 4\end{pmatrix}$
\item $\begin{pmatrix}-1 & 0\\ 4 & 3\end{pmatrix}$
\item $\begin{pmatrix}2 & 4\\ 0 & 4\end{pmatrix}$
\item $\begin{pmatrix}1 & -1\\ 0 & 2\end{pmatrix}$
\item $\begin{pmatrix}-4 & 2\\ -1 & -1\end{pmatrix}$
\item $\begin{pmatrix}0 & -8\\ 0 & 4\end{pmatrix}$
\item $\begin{pmatrix}-8 & 5\\ -10 & 7\end{pmatrix}$
\item $\begin{pmatrix}-1 & 0\\ 0 & -3\end{pmatrix}$
\item $\begin{pmatrix}-4 & 0\\ 0 & -1\end{pmatrix}$
\end{enumerate}
\end{multicols}


\emph{2. Is the matrix diagonalizable? Find its eigenvalues and state your conclusion.}
\bigskip
\begin{enumerate}[label=\arabic*., leftmargin=*]
\item $\begin{pmatrix}-2 & -6 & 3\\ 0 & 2 & 0\\ 0 & 2 & 1\end{pmatrix}$
\item $\begin{pmatrix}-3 & 0 & 0\\ 0 & 4 & 0\\ 4 & -6 & 1\end{pmatrix}$
\item $\begin{pmatrix}-2 & -4 & 0\\ 3 & 5 & 1\\ -1 & -1 & -1\end{pmatrix}$
\item $\begin{pmatrix}1 & -1 & 2\\ 4 & 2 & 2\\ -4 & 1 & -5\end{pmatrix}$
\item $\begin{pmatrix}4 & 2 & 0\\ -1 & 1 & 0\\ 2 & 2 & 2\end{pmatrix}$
\item $\begin{pmatrix}9 & 6 & -6\\ -4 & -2 & 5\\ 8 & 7 & -4\end{pmatrix}$
\item $\begin{pmatrix}-1 & 4 & -1\\ 0 & -3 & -7\\ 0 & 0 & 4\end{pmatrix}$
\item $\begin{pmatrix}4 & -5 & 5\\ 0 & -6 & 10\\ 0 & -5 & 9\end{pmatrix}$
\item $\begin{pmatrix}1 & 1 & 2\\ -10 & -6 & -8\\ 2 & 1 & 0\end{pmatrix}$
\item $\begin{pmatrix}3 & 0 & 0\\ -3 & 6 & -3\\ -6 & 6 & -3\end{pmatrix}$
\item $\begin{pmatrix}8 & -4 & 4\\ 2 & -2 & 0\\ -10 & 4 & -6\end{pmatrix}$
\item $\begin{pmatrix}-3 & 0 & 0\\ 2 & -1 & -1\\ 3 & 2 & -4\end{pmatrix}$
\item $\begin{pmatrix}-3 & 0 & 0\\ 1 & -2 & 0\\ -9 & -4 & 2\end{pmatrix}$
\item $\begin{pmatrix}-2 & -10 & 8\\ 0 & 3 & -4\\ 0 & 0 & -1\end{pmatrix}$
\item $\begin{pmatrix}-1 & 1 & -2\\ -1 & 1 & -1\\ 0 & 0 & 1\end{pmatrix}$
\item $\begin{pmatrix}4 & 0 & 0\\ 10 & -1 & 0\\ -10 & 0 & -1\end{pmatrix}$
\item $\begin{pmatrix}0 & 0 & 0\\ 1 & -1 & 0\\ 1 & 0 & -1\end{pmatrix}$
\item $\begin{pmatrix}-2 & 0 & -1\\ 0 & -2 & 0\\ 0 & 1 & -3\end{pmatrix}$
\item $\begin{pmatrix}1 & 2 & 2\\ 1 & 0 & 2\\ -2 & -2 & -4\end{pmatrix}$
\item $\begin{pmatrix}-3 & -4 & -8\\ 0 & -1 & 4\\ 0 & 0 & -3\end{pmatrix}$
\item $\begin{pmatrix}2 & 0 & 4\\ -1 & 2 & -9\\ 0 & 0 & -2\end{pmatrix}$
\item $\begin{pmatrix}3 & 0 & 0\\ -3 & -2 & -1\\ 12 & 0 & -3\end{pmatrix}$
\item $\begin{pmatrix}3 & 0 & 0\\ -6 & 0 & 0\\ -12 & 0 & -3\end{pmatrix}$
\item $\begin{pmatrix}-1 & 0 & -2\\ 0 & -3 & -1\\ 1 & 0 & -4\end{pmatrix}$
\item $\begin{pmatrix}-1 & 0 & 0\\ -1 & -2 & 0\\ -6 & -6 & 4\end{pmatrix}$
\item $\begin{pmatrix}-3 & 3 & 0\\ 0 & 0 & 0\\ 0 & 6 & 3\end{pmatrix}$
\item $\begin{pmatrix}-2 & 0 & 1\\ -1 & -1 & 2\\ 0 & 0 & -2\end{pmatrix}$
\item $\begin{pmatrix}2 & 0 & -1\\ 4 & 0 & -2\\ 2 & 0 & -1\end{pmatrix}$
\item $\begin{pmatrix}4 & -3 & 5\\ -4 & 2 & -4\\ -4 & 3 & -5\end{pmatrix}$
\item $\begin{pmatrix}5 & -1 & -2\\ 0 & 3 & 0\\ 2 & -1 & 1\end{pmatrix}$
\end{enumerate}


\emph{3. Further problems.}
\bigskip
\begin{enumerate}[label=\arabic*., leftmargin=*]
\item \textbf{(Powers.)} For $A=\begin{pmatrix}0 & 2\\ -1 & 3\end{pmatrix}$ find $A^{4}$ by diagonalization.
\item \textbf{(Powers.)} For $A=\begin{pmatrix}-1 & 0\\ 0 & 2\end{pmatrix}$ find $A^{3}$ by diagonalization.
\item \textbf{(Powers.)} For $A=\begin{pmatrix}2 & 0\\ -1 & 1\end{pmatrix}$ find $A^{4}$ by diagonalization.
\item \textbf{(Powers.)} For $A=\begin{pmatrix}3 & -8\\ 0 & -1\end{pmatrix}$ find $A^{4}$ by diagonalization.
\item \textbf{(Powers.)} For $A=\begin{pmatrix}2 & 0\\ 0 & 1\end{pmatrix}$ find $A^{4}$ by diagonalization.
\item \textbf{(Powers.)} For $A=\begin{pmatrix}7 & -4\\ 8 & -5\end{pmatrix}$ find $A^{5}$ by diagonalization.
\item \textbf{(Fibonacci numbers.)} Find the eigenvalues of $F=\left(\begin{smallmatrix}1&1\\1&0\end{smallmatrix}\right)$ and explain why the entries of $F^n$ are Fibonacci numbers.
\item \textbf{(A non-diagonalizable matrix.)} Why is $\left(\begin{smallmatrix}2&1\\0&2\end{smallmatrix}\right)$ not diagonalizable?
\item \textbf{(Symmetric matrix.)} Reduce $A=\left(\begin{smallmatrix}3&1\\1&3\end{smallmatrix}\right)$ to diagonal form by an orthogonal matrix.
\item \textbf{(Trace and determinant.)} Verify the relations $\operatorname{tr}A=\lambda_1+\dots+\lambda_n$ and $\det A=\lambda_1\cdots\lambda_n$ for the diagonalizable matrix $\left(\begin{smallmatrix}4&1\\2&3\end{smallmatrix}\right)$.
\end{enumerate}

